\subsection{What do published papers actually report?}
\label{sec:res_audit}

Sections~\ref{sec:res_cohort}--\ref{sec:res_null} show that reference
composition changes the reported number. A fair objection is that we
constructed a deliberately poor selection rule and then demonstrated that it is
poor. Whether the finding matters depends on an empirical question about the
field: do published papers report enough about their reference set for a reader
to know which rule was used?

We audited ten papers that evaluate medical image generation, spanning brain
MRI, chest X-ray, cardiac MRI and CT, from 2022 to 2025, and including two
papers whose subject is evaluation methodology rather than generation. For each
we recorded five properties answerable from the text: whether the size of the
real reference set is stated; whether the number of distinct subjects behind it
is stated; whether the selection procedure is described; whether train/test
splits are stated to be subject-level; and whether the reference is drawn from
the generator's training data. Table~\ref{tab:audit} gives the result.

\begin{table}[t]
\centering
\caption{Reference-set reporting in ten medical image generation papers.
\cmark~= stated, $\sim$~= partially stated, \xmark~= not stated,
``?''~= cannot be determined from the text. One paper uses no distributional
metric and is marked --- in the metric-specific columns. Percentages exclude
that paper.}
\label{tab:audit}
\footnotesize
\begin{tabular}{llccccc}
\toprule
Paper & Metric & $N$ imgs & $N$ subj & Selection & Subj.\ split & Ref.\ from train \\
\midrule
Pinaya et al.\ 2022        & FID       & \xmark & \xmark & \xmark & \xmark & ? \\
Khader et al.\ 2023        & MS-SSIM   & ---    & ---    & ---    & \xmark & --- \\
Akbar et al.\ 2023         & FID       & \xmark & $\sim$ & \xmark & \cmark & ? \\
Kong et al.\ 2023          & FID       & \xmark & \xmark & \xmark & \xmark & ? \\
Woodland et al.\ 2024      & Fr\'echet & $\sim$ & \xmark & $\sim$ & \xmark & ? \\
Hashmi et al.\ 2024        & FID       & \xmark & \xmark & \xmark & \xmark & $\sim$ \\
BrainMRDiff 2025           & MMD       & \xmark & \xmark & \xmark & \cmark & ? \\
FADM 2024                  & FID       & \xmark & $\sim$ & \xmark & \cmark & ? \\
Rajesh et al.\ 2025        & FID, ECS  & \xmark & \xmark & \xmark & \xmark & \cmark \\
Perceptual Eval.\ 2025     & reader    & $\sim$ & \xmark & \xmark & \xmark & ? \\
\midrule
\multicolumn{2}{l}{Fully stated}       & $0/9$  & $0/9$  & $0/9$  & $3/10$ & $1/9$ \\
\bottomrule
\end{tabular}
\end{table}

Three results stand out. \emph{No paper in the sample states how many distinct
subjects its real reference set is drawn from.} No paper fully states the
reference set size, and none describes the selection procedure well enough to
reproduce it. Only three of ten state that their splits are subject-level, and
only one states explicitly whether the reference overlaps the generator's
training data --- and that one does so because the overlap is a deliberate part
of its design.

The consequence is not that these papers are wrong. It is that a reader cannot
distinguish a reference drawn from $17$ subjects from one drawn from $400$,
and Section~\ref{sec:res_cohort} shows those are different instruments giving
answers that differ by a factor of two. The information needed to tell them
apart is, in this sample, never present. That is what makes reference
composition a reporting problem rather than a modelling one, and it is why our
recommendations in Section~\ref{sec:protocol} are about what to report rather
than about what to compute.

\paragraph{Scope of the audit}
This is a purposive sample of ten papers assembled from literature searches on
medical image synthesis with distributional metrics, not a systematic review
with a pre-registered protocol, and it is small enough that the percentages
should be read as indicative rather than as prevalence estimates. Judgements
are made from the main text; a detail buried in a supplement or a code release
would be scored \xmark\ here. We report the per-paper evidence for every
judgement in the released audit record so that individual scores can be
checked and disputed.
